% Part: computability
% Chapter: recursive-functions
% Section: bounded-minimization

\documentclass[../../../include/open-logic-section]{subfiles}

\begin{document}

\olfileid[fr]{cmp}{rec}{bmi}
\olsection{Minimisation bornée}

\begin{explain}
Il est souvent utile de définir une fonction par le plus petit
nombre qui satisfait une propriété ou une relation~$P$. Si
$P$ est décidable, on peut chercher ce nombre en testant
successivement $0$, $1$, $2$, $\dots$, jusqu'à trouver le
premier qui satisfait~$P$. Sans témoin, cette recherche ne termine pas ;
elle ne garantit pas, en général, une fonction récursive primitive.
Mais si l'on ne cherche que le plus petit nombre
\emph{inférieur à une borne donnée indépendamment}, on reste
dans ce cadre. Plus généralement, supposons que $R(x,z)$
soit une relation récursive primitive. Considérons la
fonction qui associe à $x$ et~$y$ le plus petit $z<y$
tel que $R(x,z)$. On peut la calculer en testant
$R(x,0)$, $R(x,1)$, $\dots$, $R(x,y-1)$.
Pourquoi est-elle elle-même récursive primitive ?
\end{explain}

\begin{prop}
Si $R(\vec x,z)$ est récursive primitive, il en va de même
de la fonction $m_R(\vec{x},y)$ qui renvoie le plus petit
$z<y$ tel que $R(\vec x,z)$, s'il en existe un, et
la borne~$y$ sinon. Nous noterons cette fonction~$m_R$ :
\[
\bmin{z < y}{R(\vec{x}, z)},
\]
\end{prop}

\begin{proof}
Il n'existe aucun $z<0$; en particulier, aucun tel
$z$ ne satisfait $R(\vec x,z)$. On a donc
$m_R(\vec x,0)=0$.

Pour une borne de la forme $y+1$, distinguons trois cas :
\begin{enumerate}
\item Il existe $z<y$ tel que $R(\vec x,z)$.
Alors $m_R(\vec x,y+1)=m_R(\vec x,y)$.
\item Aucun $z<y$ ne convient, mais $R(\vec x,y)$
vaut. Alors $m_R(\vec x,y+1)=y$.
\item Aucun $z<y+1$ ne convient. Alors
$m_R(\vec x,y+1)=y+1$.\footnote{Correction éditoriale : le tuple de paramètres reste $\vec x$, là où ce troisième cas de la source écrit $\vec z$ sans le définir.}
\end{enumerate}
On peut donc définir $m_R$ par récurrence primitive :
\begin{align*}
m_R(\vec{x}, 0) & = 0\\
m_R(\vec{x}, y+1) & =
\begin{cases}
m_R(\vec{x}, y) & \text{si $m_R(\vec x,y)\neq y$}\\
y & \text{si $m_R(\vec x,y)=y$ et $R(\vec{x},y)$}\\
y+1 & \text{sinon.}
\end{cases}
\end{align*}
Notons qu'il existe un $z<y$ tel que $R(\vec x,z)$
si et seulement si $m_R(\vec x,y)\neq y$.
\end{proof}

\begin{prob}
Supposons $R(\vec x,z)$ récursive primitive. Définir
par récurrence primitive, à partir de~$\Char{R}$,
la fonction $m'_R(\vec{x},y)$ qui renvoie le plus
petit $z<y$ tel que $R(\vec x,z)$, s'il en existe un,
et~$0$ sinon.
\end{prob}

\end{document}
